% GENERATED FILE. Edit canonical lessons, then run npm run generate:books.
% canonical-source-hash: fnv1a64:c1da11dc627f1c2a
% canonical-lessons: MW-C19-hear-roti, MW-W19-tta, MW-C19-roti, MW-C19-hear-baati, MW-C19-baati, MW-R19-food-new-two, MW-C19-food-five, MW-R19-script-close

\chapter{Bread and Rolls: One New Consonant}
\label{ch:mw-food-five}
\hlchaptermodality{24}

\begin{chapteropening}
\textbf{By the end of this chapter:} \emph{I can add flatbread and hard wheat rolls to the food map after learning only \mw{ट} and reusing every other sign.}

Two more source-attested food words begin by ear, add only \mw{ट}, and join the earlier three in a non-compensatory four-skill payoff.
\end{chapteropening}

\section[roṭī]{Hear \emph{roṭī} — flatbread}

\label{lesson:MW-C19-hear-roti}



\begin{quote}
Recall the three-word food payoff, then say hand and lentils.
\end{quote}

\subsection*{What to know first}
\begin{quote}
\emph{roṭī} — \textbf{flatbread}
\end{quote}

The middle consonant is made with the tongue tip curled back, further back than the \emph{d} of \emph{dāl}. Hear it twice and say it once. The spelling waits for the next lesson, because one sign in it has not been taught yet.

\begin{tcolorbox}[breakable,colback=teal!4,colframe=teal!35!black,title={Before you move on}]
Source: \href{https://www.marwaripathshala.com/marwari-lesson-6-english}{Marwari Pathshala, Lesson 6}.
\end{tcolorbox}
\section[ṭa]{\mw{ट} — the curled \emph{t} in flatbread}

\label{lesson:MW-W19-tta}



\begin{quote}
Say flatbread, recall the three-word transport payoff, then write hand.
\end{quote}

\subsection*{Script — one sign}
\begin{quote}
\textbf{\mw{ट}} — \emph{ṭa}
\end{quote}

A single closed loop hanging under the headline. Set it beside the \textbf{\mw{ठ}} of \textbf{\mw{ठेला}}: \textbf{\mw{ठ}} carries a second inner ring and a puff of breath, \textbf{\mw{ट}} has neither.

\subsection*{Your turn}
Trace \textbf{\mw{ट}} once. Copy it once. Write \textbf{\mw{ठ}} beside it and say which one has the extra ring. Cover both, wait five seconds, and write \textbf{\mw{ट}} from the curled sound cue.

\begin{tcolorbox}[breakable,colback=teal!4,colframe=teal!35!black,title={Before you move on}]
Source: \href{https://www.unicode.org/charts/PDF/U0900.pdf}{Unicode Devanagari chart}.
\end{tcolorbox}
\section[roṭī]{\mw{रोटी} — flatbread on the page}

\label{lesson:MW-C19-roti}



\begin{quote}
Say flatbread and money, then write \textbf{\mw{ट}}.
\end{quote}

\subsection*{What to know first}
\begin{quote}
\textbf{\mw{रो}} + \textbf{\mw{टी}} $\to$ \textbf{\mw{रोटी}} — \emph{roṭī} — flatbread
\end{quote}

The \textbf{\mw{ो}} on \textbf{\mw{र}} and the \textbf{\mw{ी}} on \textbf{\mw{ट}} are both long vowels met earlier.

\subsection*{Your turn}
Look, cover, wait five seconds, and write \textbf{\mw{रोटी}}. Check that the second consonant has no inner ring.

\begin{tcolorbox}[breakable,colback=teal!4,colframe=teal!35!black,title={Before you move on}]
Source: \href{https://www.marwaripathshala.com/marwari-lesson-6-english}{Marwari Pathshala, Lesson 6}.
\end{tcolorbox}
\section[bāṭī]{Hear \emph{bāṭī} — hard wheat rolls}

\label{lesson:MW-C19-hear-baati}



\begin{quote}
Write flatbread and money, then say rickshaw.
\end{quote}

\subsection*{What to know first}
\begin{quote}
\emph{bāṭī} — \textbf{hard wheat rolls}
\end{quote}

Two beats, both long, with the same curled middle consonant heard in \emph{roṭī}. These are baked hard rather than soft, which is why they take a different name from ordinary flatbread. Hear it twice and say it once.

\begin{tcolorbox}[breakable,colback=teal!4,colframe=teal!35!black,title={Before you move on}]
Source: \href{https://www.marwaripathshala.com/marwari-lesson-6-english}{Marwari Pathshala, Lesson 6}.
\end{tcolorbox}
\section[bāṭī]{\mw{बाटी} — hard wheat rolls on the page}

\label{lesson:MW-C19-baati}



\begin{quote}
Say hard wheat rolls, then write flatbread and \textbf{\mw{श}}.
\end{quote}

\subsection*{What to know first}
\begin{quote}
\textbf{\mw{बा}} + \textbf{\mw{टी}} $\to$ \textbf{\mw{बाटी}} — \emph{bāṭī} — hard wheat rolls
\end{quote}

Compare \textbf{\mw{रोटी}} and \textbf{\mw{बाटी}} on one line: the last two units are identical and only the opening syllable differs.

\subsection*{Your turn}
Look, cover, wait five seconds, and write \textbf{\mw{बाटी}}. Then write \textbf{\mw{रोटी}} beneath it and check that only the first syllable changed.

\begin{tcolorbox}[breakable,colback=teal!4,colframe=teal!35!black,title={Before you move on}]
Source: \href{https://www.marwaripathshala.com/marwari-lesson-6-english}{Marwari Pathshala, Lesson 6}.
\end{tcolorbox}
\section[Practice]{Retrieve flatbread and hard wheat rolls}

\label{lesson:MW-R19-food-new-two}



\begin{quote}
Recall the five-word travel payoff, then write \textbf{\mw{ट}} and rickshaw.
\end{quote}

\subsection*{Your turn}
Hear both words in both orders, give each meaning, read two cards, then write both from sound.

\begin{tcolorbox}[breakable,colback=teal!4,colframe=teal!35!black,title={Before you move on}]
Source: \href{https://www.marwaripathshala.com/marwari-lesson-6-english}{Marwari Pathshala, Lesson 6}.
\end{tcolorbox}
\section[Practice]{Payoff — five food words in four skills}

\label{lesson:MW-C19-food-five}



\begin{quote}
Recall the earlier three-word food payoff, then say horse.
\end{quote}

\subsection*{Your turn}
\begin{enumerate}
\raggedright
  \item Identify five heard words.
  \item Produce five words from meaning cues.
  \item Match five printed cards to meanings.
  \item Write all five heard words without a model.
\end{enumerate}

Pass each skill separately. This vocabulary foundation does not yet claim an order in a kitchen, a request at a table, or a description of a meal.

\begin{tcolorbox}[breakable,colback=teal!4,colframe=teal!35!black,title={Before you move on}]
Source: \href{https://www.marwaripathshala.com/marwari-lesson-6-english}{Marwari Pathshala, Lesson 6}.
\end{tcolorbox}
\section[Practice]{Close the writing loop for \mw{ट}}

\label{lesson:MW-R19-script-close}



\begin{quote}
Recall the five-word food payoff, say wind, then write horse.
\end{quote}

\subsection*{Your turn}
Without a model, write \textbf{\mw{ट}} and \textbf{\mw{ठ}} from sound cues. Then write \textbf{\mw{रोटी}}, \textbf{\mw{बाटी}}, and \textbf{\mw{ठेला}} from dictation. Repair only the missed unit and repeat that word once.

\begin{tcolorbox}[breakable,colback=teal!4,colframe=teal!35!black,title={Before you move on}]
Sources: \href{https://www.marwaripathshala.com/marwari-lesson-6-english}{Marwari Pathshala, Lesson 6}; \href{https://www.unicode.org/charts/PDF/U0900.pdf}{Unicode Devanagari chart}.
\end{tcolorbox}
